\documentclass{article}

% to compile a preprint version, e.g., for submission to arXiv, add add the
% [preprint] option:
\usepackage[preprint]{neurips_2024}

\usepackage[utf8]{inputenc} % allow utf-8 input
\usepackage[T1]{fontenc}    % use 8-bit T1 fonts
\usepackage{hyperref}       % hyperlinks
\usepackage{url}            % simple URL typesetting
\usepackage{booktabs}       % professional-quality tables
\usepackage{amsfonts}       % blackboard math symbols
\usepackage{nicefrac}       % compact symbols for 1/2, etc.
\usepackage{microtype}      % microtypography
\usepackage{xcolor}         % colors


\title{Deep Learning Project Proposal Template}

\author{%
  Team Leader \\
  \texttt{2023001001} \\
  \And
  Team Member 1 \\
  \texttt{2023002002} \\
  \And
  Team Member 2 \\
  \texttt{2023003003} \\
}

\begin{document}
\maketitle

\begin{abstract}
  In one sentence, summarize what you want to do and what you will demonstrate.
\end{abstract}

\section{Introduction}
Introduce the context and motivation behind your project. Describe the problem you are addressing, and provide any necessary background information or related work \textit{with citations}.

\section{Goal of Achievement}
State the specific goals you aim to achieve with your project, focusing on the demonstration you plan to create and the effects you intend to showcase. Be clear and concise, and ensure your goals are measurable and achievable within the scope of your project.

\section{Methods}
Outline the methodology you plan to use to achieve your objectives. This may include a description of the algorithms, techniques, or tools you will employ. If your project involves data, provide details on the datasets you will use or generate, as well as any pre-processing steps. This is only an initial version of your method, and you can update it during the project.

\section{Timeline}
Provide an estimated timeline for your project, including the start and end dates, and any intermediate milestones. Make sure the timeline is realistic and accounts for potential challenges or delays. Don't forget that you will also need some time to prepare a poster for the final poster presentation in the last week.

\section{Workload Distribution}
In this section, provide a clear description of the workload distribution among your team members. Outline the specific tasks and responsibilities assigned to each team member, ensuring that the workload is fairly distributed and that each member contributes to the project's success.

\bibliographystyle{unsrt}
\bibliography{ref}

\end{document}
